% Options for packages loaded elsewhere
\PassOptionsToPackage{unicode}{hyperref}
\PassOptionsToPackage{hyphens}{url}
\PassOptionsToPackage{dvipsnames,svgnames,x11names}{xcolor}
\documentclass[
]{article}
\usepackage{xcolor}
\usepackage[margin=1in]{geometry}
\usepackage{amsmath,amssymb}
\setcounter{secnumdepth}{-\maxdimen} % remove section numbering
\usepackage{iftex}
\ifPDFTeX
  \usepackage[T1]{fontenc}
  \usepackage[utf8]{inputenc}
  \usepackage{textcomp} % provide euro and other symbols
\else % if luatex or xetex
  \usepackage{unicode-math} % this also loads fontspec
  \defaultfontfeatures{Scale=MatchLowercase}
  \defaultfontfeatures[\rmfamily]{Ligatures=TeX,Scale=1}
\fi
\usepackage{lmodern}
\ifPDFTeX\else
  % xetex/luatex font selection
  \setmainfont[]{Helvetica Neue}
\fi
% Use upquote if available, for straight quotes in verbatim environments
\IfFileExists{upquote.sty}{\usepackage{upquote}}{}
\IfFileExists{microtype.sty}{% use microtype if available
  \usepackage[]{microtype}
  \UseMicrotypeSet[protrusion]{basicmath} % disable protrusion for tt fonts
}{}
\makeatletter
\@ifundefined{KOMAClassName}{% if non-KOMA class
  \IfFileExists{parskip.sty}{%
    \usepackage{parskip}
  }{% else
    \setlength{\parindent}{0pt}
    \setlength{\parskip}{6pt plus 2pt minus 1pt}}
}{% if KOMA class
  \KOMAoptions{parskip=half}}
\makeatother
\usepackage{longtable,booktabs,array}
\usepackage{caption}
\captionsetup[table]{skip=6pt}
\newcounter{none} % for unnumbered tables
\usepackage{calc} % for calculating minipage widths
% Correct order of tables after \paragraph or \subparagraph
\usepackage{etoolbox}
\makeatletter
\patchcmd\longtable{\par}{\if@noskipsec\mbox{}\fi\par}{}{}
\makeatother
\usepackage{graphicx}
\makeatletter
\newsavebox\pandoc@box
\newcommand*\pandocbounded[1]{% scales image to fit in text height/width
  \sbox\pandoc@box{#1}%
  \Gscale@div\@tempa{\textheight}{\dimexpr\ht\pandoc@box+\dp\pandoc@box\relax}%
  \Gscale@div\@tempb{\linewidth}{\wd\pandoc@box}%
  \ifdim\@tempb\p@<\@tempa\p@\let\@tempa\@tempb\fi% select the smaller of both
  \ifdim\@tempa\p@<\p@\scalebox{\@tempa}{\usebox\pandoc@box}%
  \else\usebox{\pandoc@box}%
  \fi%
}
% Set default figure placement to htbp
\def\fps@figure{htbp}
\makeatother
\setlength{\emergencystretch}{3em} % prevent overfull lines
\providecommand{\tightlist}{%
  \setlength{\itemsep}{0pt}\setlength{\parskip}{0pt}}
\usepackage{bookmark}
\IfFileExists{xurl.sty}{\usepackage{xurl}}{} % add URL line breaks if available
\urlstyle{same}
% fallback for those not using the hyperref driver hyperxmp:
\makeatletter
\@ifundefined{xmpquote}{\newcommand{\xmpquote}[1]{#1}}{}
\makeatother
\hypersetup{
  pdftitle={Reporte de Avance No. 2},
  pdfauthor={Equipo de Desarrollo ABAQ},
  colorlinks=true,
  linkcolor={Maroon},
  filecolor={Maroon},
  citecolor={Blue},
  urlcolor={Blue},
  pdfcreator={LaTeX via pandoc}}

\title{Reporte de Avance No. 2}
\author{Equipo de Desarrollo ABAQ}
\date{}

\begin{document}
\maketitle

{
\setcounter{tocdepth}{3}
\tableofcontents
}
\section{Reporte de Avance No.~2}\label{reporte-de-avance-no.-2}

{\def\LTcaptype{none} % do not increment counter
\begin{longtable}[]{@{}ll@{}}
\toprule\noalign{}
\endhead
\bottomrule\noalign{}
\endlastfoot
\textbf{Fecha} & 9 de septiembre de 2026 \\
\textbf{Versión} & v1.1.0 \\
\textbf{Estado} & En desarrollo (rama \texttt{dev}) \\
\textbf{Preparado por} & Equipo de Desarrollo ABAQ \\
\end{longtable}
}

\begin{center}\rule{0.5\linewidth}{0.5pt}\end{center}

\textbf{Descargas} Puedes descargar este documento en los siguientes
formatos:

\begin{verbatim}
-  **[Descargar PDF](latex/report2.pdf)**
-  **[Descargar Word (.docx)](latex/report2.docx)**
\end{verbatim}

\subsection{Resumen ejecutivo}\label{resumen-ejecutivo}

Este documento presenta el segundo reporte de avance del proyecto ABAQ,
la plataforma digital para la gestión del servicio social. Durante este
período se realizaron mejoras significativas tanto en la seguridad del
servidor como en la experiencia visual de la aplicación web.

En el lado del servidor se implementó un sistema de autenticación que
protege toda la información sensible: ahora solo los usuarios con sesión
iniciada pueden consultar o modificar datos de estudiantes, contactos,
actividades y encuestas. La única excepción intencional es el formulario
público para aspirantes, que continúa disponible sin restricción.

En la aplicación web se construyeron tres pantallas completamente
nuevas: el \textbf{Panel de Administración}, la \textbf{Vista de Detalle
de Estudiante} y la \textbf{Vista de Detalle de Actividad}. Además se
mejoraron las listas existentes, se rediseñó la página de inicio y se
integró la capacidad de subir documentos directamente desde la
plataforma.

El presente reporte documenta y evidencia estos avances con capturas de
pantalla del sistema en funcionamiento, cumpliendo con el compromiso de
coordinación y seguimiento del progreso establecido con el cliente.

\begin{center}\rule{0.5\linewidth}{0.5pt}\end{center}

\subsection{1. Introducción}\label{introducciuxf3n}

Este documento resume los cambios más recientes realizados en la
plataforma \textbf{ABAQ}, la herramienta de gestión de servicio social.
El objetivo es comunicar, de forma clara y sin tecnicismos, qué mejoró,
qué se agregó y qué quedó más seguro desde la última entrega.

Los cambios se dividen en dos áreas principales:

\begin{itemize}
\tightlist
\item
  \textbf{Servidor (Backend):} la parte que no se ve, encargada de
  guardar y proteger los datos.
\item
  \textbf{Aplicación web (Frontend):} la parte visual con la que
  interactúan administradores y estudiantes.
\end{itemize}

\begin{center}\rule{0.5\linewidth}{0.5pt}\end{center}

\subsection{2. Cambios en el Servidor
(Backend)}\label{cambios-en-el-servidor-backend}

\subsubsection{Protección de rutas con
autenticación}\label{protecciuxf3n-de-rutas-con-autenticaciuxf3n}

Anteriormente, alguien con conocimientos técnicos podía acceder a la
información de estudiantes, contactos, encuestas y actividades sin
necesidad de haber iniciado sesión. Ahora eso ya no es posible.

Se protegieron los siguientes módulos:

\begin{itemize}
\tightlist
\item
  \textbf{Estudiantes:} para ver, crear, editar o eliminar un
  estudiante, es obligatorio haber iniciado sesión.
\item
  \textbf{Contactos:} toda acción sobre contactos requiere
  autenticación.
\item
  \textbf{Encuestas:} la consulta de encuestas ya guardadas ahora
  requiere sesión iniciada. La encuesta pública que llenan los
  aspirantes sigue disponible sin restricción.
\item
  \textbf{Actividades y mascotas:} protegidas de la misma manera.
\end{itemize}

\subsubsection{Mejoras en la
configuración}\label{mejoras-en-la-configuraciuxf3n}

\begin{itemize}
\tightlist
\item
  Se actualizó el archivo de ejemplo de configuración
  (\texttt{.env.example}) para incluir las variables necesarias para el
  sistema de autenticación JWT.
\item
  Se actualizó el \texttt{README} con instrucciones más claras sobre
  cómo instalar y configurar el servidor.
\item
  Se excluyó un archivo de registro local (\texttt{log.txt}) del control
  de versiones.
\end{itemize}

\subsubsection{Resumen de cambios ---
Backend}\label{resumen-de-cambios-backend}

{\def\LTcaptype{none} % do not increment counter
\begin{longtable}[]{@{}
  >{\raggedright\arraybackslash}p{(\linewidth - 2\tabcolsep) * \real{0.5000}}
  >{\raggedright\arraybackslash}p{(\linewidth - 2\tabcolsep) * \real{0.5000}}@{}}
\toprule\noalign{}
\begin{minipage}[b]{\linewidth}\raggedright
Área
\end{minipage} & \begin{minipage}[b]{\linewidth}\raggedright
Descripción del cambio
\end{minipage} \\
\midrule\noalign{}
\endhead
\bottomrule\noalign{}
\endlastfoot
Seguridad & Se requiere inicio de sesión para acceder a datos de
estudiantes, contactos, encuestas y actividades \\
Encuesta pública & El formulario para aspirantes continúa accesible sin
sesión \\
Documentación & README actualizado con instrucciones de instalación y
variables de entorno \\
Configuración & Archivo \texttt{.env.example} con variables de
autenticación JWT \\
Limpieza & Archivo \texttt{log.txt} excluido del repositorio \\
\end{longtable}
}

\begin{center}\rule{0.5\linewidth}{0.5pt}\end{center}

\subsection{3. Cambios en la Aplicación Web
(Frontend)}\label{cambios-en-la-aplicaciuxf3n-web-frontend}

\subsubsection{Panel de
Administración}\label{panel-de-administraciuxf3n}

Se creó una pantalla de inicio exclusiva para administradores. Muestra
métricas generales (estudiantes, actividades, inscripciones,
asignaciones pendientes), una agenda con próximas actividades y un panel
de seguimiento con elementos que requieren atención.

\subsubsection{Vista de Detalle de
Estudiante}\label{vista-de-detalle-de-estudiante}

Nueva pantalla con el expediente completo del estudiante: datos
personales, horas requeridas, horas acreditadas, edición de información,
carga de documentos oficiales (carta de aceptación, constancia de
plática, carta de término) y listado de actividades.

\subsubsection{Lista de Estudiantes}\label{lista-de-estudiantes}

Tabla mejorada con buscador por nombre, carrera, correo o teléfono, y
acceso directo al detalle de cada estudiante.

\subsubsection{Lista de Actividades}\label{lista-de-actividades}

Tarjetas mejoradas con buscador por título, descripción, tipo, horas o
ID. Cada actividad muestra tipo, duración, horario y enlace.

\subsubsection{Panel del Estudiante}\label{panel-del-estudiante}

Rediseño de la pantalla de inicio para estudiantes: saludo
personalizado, carrera, semestre, horas requeridas, horas acreditadas,
actividades pendientes, historial y gráficas de progreso.

\subsubsection{Carga de Archivos}\label{carga-de-archivos}

Nuevo servicio para subir documentos PDF directamente desde la
plataforma sin salir de la aplicación.

\subsubsection{Resumen de cambios ---
Frontend}\label{resumen-de-cambios-frontend}

{\def\LTcaptype{none} % do not increment counter
\begin{longtable}[]{@{}
  >{\raggedright\arraybackslash}p{(\linewidth - 2\tabcolsep) * \real{0.5000}}
  >{\raggedright\arraybackslash}p{(\linewidth - 2\tabcolsep) * \real{0.5000}}@{}}
\toprule\noalign{}
\begin{minipage}[b]{\linewidth}\raggedright
Área
\end{minipage} & \begin{minipage}[b]{\linewidth}\raggedright
Descripción del cambio
\end{minipage} \\
\midrule\noalign{}
\endhead
\bottomrule\noalign{}
\endlastfoot
Panel de administración & Nueva pantalla con métricas, agenda y panel de
seguimiento \\
Detalle de estudiante & Nueva pantalla con expediente, edición y carga
de documentos \\
Lista de estudiantes & Tabla mejorada con buscador y acceso a detalle \\
Lista de actividades & Tarjetas mejoradas con buscador por múltiples
criterios \\
Panel del estudiante & Rediseño con datos personales, progreso e
historial \\
Carga de archivos & Nuevo servicio para subir documentos PDF \\
Navegación & Nuevas rutas para las pantallas de detalle \\
\end{longtable}
}

\begin{center}\rule{0.5\linewidth}{0.5pt}\end{center}

\subsection{4. Evidencia Visual}\label{evidencia-visual}

\subsubsection{Panel de
administración}\label{panel-de-administraciuxf3n-1}

\begin{figure}
\centering
\pandocbounded{\includegraphics[keepaspectratio,alt={Panel de administración --- métricas generales}]{attachments/Capture-2026-09-10-180327.png}}
\caption{Panel de administración --- métricas generales}
\end{figure}

\begin{figure}
\centering
\pandocbounded{\includegraphics[keepaspectratio,alt={Panel de administración --- agenda y resumen}]{attachments/Capture-2026-09-10-180430.png}}
\caption{Panel de administración --- agenda y resumen}
\end{figure}

\subsubsection{Gestión de estudiantes}\label{gestiuxf3n-de-estudiantes}

\begin{figure}
\centering
\pandocbounded{\includegraphics[keepaspectratio,alt={Lista de estudiantes con buscador y horas acreditadas}]{attachments/Capture-2026-09-10-180521.png}}
\caption{Lista de estudiantes con buscador y horas acreditadas}
\end{figure}

\begin{figure}
\centering
\pandocbounded{\includegraphics[keepaspectratio,alt={Expediente del estudiante --- información completa}]{attachments/Capture-2026-09-10-180604.png}}
\caption{Expediente del estudiante --- información completa}
\end{figure}

\begin{figure}
\centering
\pandocbounded{\includegraphics[keepaspectratio,alt={Expediente del estudiante --- documentos del expediente}]{attachments/Capture-2026-09-10-180639.png}}
\caption{Expediente del estudiante --- documentos del expediente}
\end{figure}

\begin{figure}
\centering
\pandocbounded{\includegraphics[keepaspectratio,alt={Expediente del estudiante --- documentos y actividades}]{attachments/Capture-2026-09-10-180830.png}}
\caption{Expediente del estudiante --- documentos y actividades}
\end{figure}

\subsubsection{Actividades disponibles}\label{actividades-disponibles}

\begin{figure}
\centering
\pandocbounded{\includegraphics[keepaspectratio,alt={Lista de actividades --- parte 1}]{attachments/Capture-2026-09-10-180920.png}}
\caption{Lista de actividades --- parte 1}
\end{figure}

\begin{figure}
\centering
\pandocbounded{\includegraphics[keepaspectratio,alt={Lista de actividades --- parte 2}]{attachments/Capture-2026-09-10-180950.png}}
\caption{Lista de actividades --- parte 2}
\end{figure}

\subsubsection{Panel del estudiante}\label{panel-del-estudiante-1}

\begin{figure}
\centering
\pandocbounded{\includegraphics[keepaspectratio,alt={Panel del estudiante --- bienvenida y perfil}]{attachments/Capture-2026-09-10-181022.png}}
\caption{Panel del estudiante --- bienvenida y perfil}
\end{figure}

\begin{figure}
\centering
\pandocbounded{\includegraphics[keepaspectratio,alt={Panel del estudiante --- actividades pendientes}]{attachments/Capture-2026-09-10-181109.png}}
\caption{Panel del estudiante --- actividades pendientes}
\end{figure}

\begin{figure}
\centering
\pandocbounded{\includegraphics[keepaspectratio,alt={Panel del estudiante --- historial de actividades}]{attachments/Capture-2026-09-10-181123.png}}
\caption{Panel del estudiante --- historial de actividades}
\end{figure}

\begin{figure}
\centering
\pandocbounded{\includegraphics[keepaspectratio,alt={Panel del estudiante --- progreso y actividades asignadas}]{attachments/Capture-2026-09-10-181135.png}}
\caption{Panel del estudiante --- progreso y actividades asignadas}
\end{figure}

\begin{center}\rule{0.5\linewidth}{0.5pt}\end{center}

\subsection{5. Estado Actual del
Proyecto}\label{estado-actual-del-proyecto}

\textbf{Sigue en discusión} El estado detallado de los módulos se
definirá en la siguiente reunión de seguimiento.

\begin{center}\rule{0.5\linewidth}{0.5pt}\end{center}

\subsection{6. Próximos Pasos}\label{pruxf3ximos-pasos}

Los cambios descritos en este reporte han sido revisados, documentados y
evidenciados con capturas de pantalla del sistema en funcionamiento. El
compromiso de coordinación y seguimiento del progreso queda cubierto con
la entrega de este documento.

\textbf{No existen pendientes adicionales para esta etapa.}

\end{document}
